\chapter{Python, PyTorch, and Deep Learning Foundations}
\label{chap:ch03}

\section{Object-Oriented Deep Learning in PyTorch}
PyTorch is a modern, dynamic scientific computing library that combines multidimensional tensor mathematics on NVIDIA CUDA GPUs with an automatic differentiation engine (\texttt{torch.autograd}).

Every neural network component in PyTorch is implemented as an Object-Oriented class that inherits from the fundamental base class: \texttt{torch.nn.Module}.

\begin{conceptbox}{Analogy: A Factory Assembly Line Worker}
Think of an \texttt{nn.Module} as a specialized worker on an assembly line:
\begin{itemize}
    \item \textbf{Internal tools (Parameters):} The worker has specific adjustable tools (weights and biases) initialized in the \texttt{\_\_init\_\_()} method.
    \item \textbf{Action (\texttt{forward()}):} When raw materials (input tensors) arrive on the conveyor belt, the worker applies their tools to transform the material and passes it to the next workstation.
    \item \textbf{Fine-tuning (Gradients):} At the end of the shift, quality control provides feedback on errors, and the worker slightly adjusts their tool settings (updating weights via backpropagation) to perform better tomorrow!
\end{itemize}
\end{conceptbox}

\section{The Anatomy of an \texttt{nn.Module}}
Here is how a clean, minimal PyTorch module is constructed:
\begin{lstlisting}[language=Python]
import torch
import torch.nn as nn

class SimpleModalityStem(nn.Module):
    def __init__(self, in_channels=2, out_channels=16):
        super(SimpleModalityStem, self).__init__()
        # Define internal trainable layers
        self.conv = nn.Conv2d(in_channels, out_channels, kernel_size=3, padding=1)
        self.relu = nn.ReLU(inplace=True)

    def forward(self, x):
        # x arrives with shape: [Batch, InChannels, Height, Width]
        feat = self.conv(x)
        out = self.relu(feat)
        return out
\end{lstlisting}

\subsection{Why Never Call \texttt{.forward()} Directly}
In Python, one calls the instantiated model like a function:
\begin{lstlisting}[language=Python]
stem = SimpleModalityStem(2, 16)
output = stem(my_input_tensor) # Correct: invokes __call__()
\end{lstlisting}
Under the hood, PyTorch's \texttt{\_\_call\_\_} method registers forward pre-hooks, execution profiling, and backward autograd hooks before calling \texttt{forward(x)}. Calling \texttt{stem.forward(x)} directly bypasses these crucial internal mechanisms!

\section{Essential Tensor Operations in Satellite Computer Vision}

\subsection{Broadcasting}
When performing arithmetic operations between tensors of different shapes, PyTorch automatically stretches (broadcasts) the smaller tensor across the missing or singleton dimensions without copying data in memory.
\begin{mathbox}{Numerical Broadcasting Example}
Suppose we have a predicted optical feature map $\mathbf{P} \in \mathbb{R}^{2 \times 4 \times 128 \times 128}$ and a single-channel cloud mask $\mathbf{M} \in \mathbb{R}^{2 \times 1 \times 128 \times 128}$.
We evaluate:
\[
\mathbf{Y} = \mathbf{M} \odot \mathbf{P}
\]
Along dimension $C=1$, $\mathbf{M}$ has 1 channel while $\mathbf{P}$ has 4 channels. PyTorch broadcasts $\mathbf{M}$ identically across all 4 channels (Red, Green, Blue, NIR):
\[
\mathbf{Y}_{b, c, h, w} = \mathbf{M}_{b, 0, h, w} \times \mathbf{P}_{b, c, h, w} \quad \text{for each } c \in \{0, 1, 2, 3\}
\]
\end{mathbox}

\subsection{Channel Concatenation (\texttt{torch.cat})}
Combining different physical modalities (such as SAR features and optical features) is achieved by stacking tensors along the channel dimension ($C=1$):
\begin{lstlisting}[language=Python]
# s1_feat: [B, 16, H, W]
# s2_feat: [B, 16, H, W]
fused = torch.cat([s1_feat, s2_feat], dim=1) # Shape: [B, 32, H, W]
\end{lstlisting}
Spatial dimensions ($H, W$) and batch size ($B$) must match exactly.

\section{The Core PyTorch Training Step}
Every deep learning training step consists of six coordinated phases:
\begin{lstlisting}[language=Python]
# 1. Zero out old gradients from the previous batch
optimizer.zero_grad()

# 2. Forward pass: compute predictions
pred = model(x_input)

# 3. Loss calculation: evaluate scalar error
loss = criterion(pred, target)

# 4. Backward pass: compute dLoss/dWeight for all parameters
loss.backward()

# 5. Gradient clipping: prevent exploding gradients
torch.nn.utils.clip_grad_norm_(model.parameters(), max_norm=1.0)

# 6. Optimizer step: update weights in the direction of negative gradient
optimizer.step()
\end{lstlisting}

\section{Summary of Chapter 3}
\begin{enumerate}
    \item Neural network building blocks in PyTorch inherit from \texttt{torch.nn.Module}.
    \item Layers and trainable weights are initialized in \texttt{\_\_init\_\_()}; execution logic resides in \texttt{forward()}.
    \item Always call \texttt{model(x)} rather than \texttt{model.forward(x)} to preserve autograd hooks.
    \item Broadcasting allows single-channel spatial masks to multiply multi-band optical tensors seamlessly.
    \item The training cycle follows: \texttt{zero\_grad()} $\to$ forward pass $\to$ loss calculation $\to$ \texttt{loss.backward()} $\to$ \texttt{optimizer.step()}.
\end{enumerate}

\section{Practice Exercises}
\begin{enumerate}
    \item What is the output shape when tensor $\mathbf{A}$ of shape $[2, 1, 64, 64]$ is multiplied element-wise with tensor $\mathbf{B}$ of shape $[2, 32, 64, 64]$?
    \item What will happen during training if you forget to call \texttt{optimizer.zero\_grad()} before \texttt{loss.backward()}?
    \item Why does \texttt{torch.cat([tensor1, tensor2], dim=1)} require that dimensions 0, 2, and 3 match identically?
\end{enumerate}

\subsection*{Solutions}
\begin{enumerate}
    \item Shape: $[2, 32, 64, 64]$. The singleton channel dimension of $\mathbf{A}$ is broadcast across the 32 channels of $\mathbf{B}$.
    \item Gradients in PyTorch accumulate by default (they are summed with existing values in \texttt{param.grad}). If \texttt{optimizer.zero\_grad()} is omitted, gradients from the current batch will add to gradients from all preceding batches, causing the effective gradient magnitude to blow up and destabilize training.
    \item Concatenation joins tensors along a specific axis without changing the geometric grid along any other axis. If height, width, or batch size differed, the elements could not form a valid rectangular multidimensional array.
\end{enumerate}
